\documentclass{article}
\usepackage{amsmath,amssymb,amsthm}
\usepackage{algorithm}
\usepackage{algorithmic}
\usepackage{xcolor}
\newtheorem{theorem}{Theorem}
\newtheorem{lemma}{Lemma}
\newcommand{\E}{\mathbb{E}}
\newcommand{\R}{\mathbb{R}}
\newcommand{\norm}[1]{\left\|#1\right\|}
\newcommand{\grad}{\nabla}
\newcommand{\sig}{\sigma}
\newcommand{\clip}{\mathcal{C}}
\begin{document}

\section*{Algorithm Name}
\textbf{Differentially Private Stochastic Gradient Descent (DP‑SGD) with Gaussian Noise}

\section*{Mathematical Formulation}
Let $f:\R^{d}\rightarrow\R$ be the empirical risk 
\[
f(w)=\frac{1}{n}\sum_{i=1}^{n}\ell(w;x_{i}),
\]
where each loss $\ell(\cdot;x_i)$ is $L$–smooth.  
At iteration $t$ a minibatch $B_t\subset\{1,\dots,n\}$ of size $b$ is sampled uniformly without replacement.  

\begin{enumerate}
\item \textbf{Clipping:} For each $i\in B_t$ compute the (per‑example) gradient
\[
g_i(w_t)=\grad\ell(w_t;x_i),
\]
and clip it to a norm bound $C>0$:
\[
\tilde g_i(w_t)=g_i(w_t)\cdot 
\min\!\Bigl\{1,\frac{C}{\norm{g_i(w_t)}}\Bigr\}.
\tag{1}
\]

\item \textbf{Averaging and noise addition:}
\[
\bar g_t=\frac{1}{b}\sum_{i\in B_t}\tilde g_i(w_t),\qquad 
z_t\sim\mathcal N(0,\sig^2 C^{2}\mathbf I_{d}),\qquad 
\tilde g_t=\bar g_t+z_t.
\tag{2}
\]

\item \textbf{Parameter update:}
\[
w_{t+1}=w_t-\eta_t\,\tilde g_t,
\tag{3}
\]
where $\eta_t>0$ is the learning rate.
\end{enumerate}

The Gaussian noise (2) guarantees $(\varepsilon,\delta)$‑differential privacy after $T$ iterations via the moments accountant; the privacy budget $\varepsilon$ is a function of $\sig$, $C$, $b$, $T$ and $\delta$ (omitted here for brevity).

\section*{Pseudocode}
\begin{algorithm}[H]
\caption{DP‑SGD with Gaussian Noise}
\begin{algorithmic}[1]
\STATE \textbf{Input:} data $\{x_i\}_{i=1}^n$, loss $\ell$, clipping norm $C$, noise scale $\sig$, batch size $b$, learning‑rate schedule $\{\eta_t\}_{t=0}^{T-1}$, total iterations $T$.
\STATE Initialize $w_0\in\R^d$.
\FOR{$t=0$ \TO $T-1$}
    \STATE Sample a minibatch $B_t$ of size $b$ uniformly at random.
    \FOR{each $i\in B_t$}
        \STATE Compute $g_i(w_t)=\grad\ell(w_t;x_i)$.
        \STATE Clip: $\tilde g_i(w_t)=g_i(w_t)\min\{1,\;C/\norm{g_i(w_t)}\}$.
    \ENDFOR
    \STATE $\bar g_t\gets \frac{1}{b}\sum_{i\in B_t}\tilde g_i(w_t)$.
    \STATE Sample $z_t\sim\mathcal N(0,\sig^2 C^{2}\mathbf I_d)$.
    \STATE $\tilde g_t\gets \bar g_t+z_t$.
    \STATE $w_{t+1}\gets w_t-\eta_t\tilde g_t$.
\ENDFOR
\STATE \textbf{Output:} $w_T$ (or the average $\frac1T\sum_{t=0}^{T-1}w_t$).
\end{algorithmic}
\end{algorithm}

\section*{Dry Run Example}
Consider the scalar quadratic $f(w)=(w-3)^2$.
\[
\grad f(w)=2(w-3),\qquad L=2.
\]
Set $C=10$ (no clipping needed), $\sig=0$ (to illustrate the deterministic part), batch size $b=1$, learning rate $\eta=0.1$, and run three iterations from $w_0=0$.

\begin{center}
\begin{tabular}{c|c|c|c}
$t$ & $g_t=\grad f(w_t)$ & $w_{t+1}=w_t-\eta g_t$ & $f(w_{t+1})$\\ \hline
0 & $2(0-3)=-6$ & $0-0.1(-6)=0.6$ & $(0.6-3)^2=5.76$\\
1 & $2(0.6-3)=-4.8$ & $0.6-0.1(-4.8)=1.08$ & $(1.08-3)^2=3.6864$\\
2 & $2(1.08-3)=-3.84$ & $1.08-0.1(-3.84)=1.464$ & $(1.464-3)^2=2.3617$\\
\end{tabular}
\end{center}
The objective value decreases at each step, illustrating the basic SGD dynamics. Adding Gaussian noise $\mathcal N(0,\sig^2 C^2)$ would perturb the updates while preserving privacy.

\newpage
\section*{Assumptions}
\begin{enumerate}
\item \textbf{Smoothness.} $f$ is $L$‑smooth:
\[
\forall u,v\in\R^d,\qquad 
f(v)\le f(u)+\langle\grad f(u),v-u\rangle+\frac{L}{2}\norm{v-u}^2.
\tag{A1}
\]

\item \textbf{Bounded variance of stochastic gradient after clipping.} Define the unbiased minibatch estimator
\[
\bar g_t=\frac{1}{b}\sum_{i\in B_t}\tilde g_i(w_t),\qquad 
\E[\bar g_t\,|\,w_t]=\grad f(w_t).
\]
Assume
\[
\E\bigl[\norm{\bar g_t-\grad f(w_t)}^2\mid w_t\bigr]\le \nu^2,
\tag{A2}
\]
for some $\nu\ge0$ (depends on $C$, $b$, and data distribution).

\item \textbf{Gaussian noise.} $z_t\sim\mathcal N(0,\sig^2 C^{2}\mathbf I_d)$ and is independent of all past iterates.

\item \textbf{Learning‑rate schedule.} Constant step size $\eta_t\equiv\eta\le \frac{1}{L}$.

\item \textbf{Existence of a global minimum.} $f^\star:=\inf_{w}f(w)>-\infty$.
\end{enumerate}

\section*{Main Convergence Theorem}
\begin{theorem}[Expected Gradient Norm for DP‑SGD]\label{thm:main}
Under Assumptions (A1)–(A5) and with a constant step size $\eta\le\frac{1}{L}$, the iterates $\{w_t\}_{t=0}^{T-1}$ generated by the algorithm satisfy
\[
\frac{1}{T}\sum_{t=0}^{T-1}\E\!\bigl[\norm{\grad f(w_t)}^2\bigr]
\;\le\;
\frac{2\bigl(f(w_0)-f^\star\bigr)}{\eta T}
\;+\;
L\eta\,\bigl(\nu^2+\sig^2 C^{2}\bigr).
\tag{4}
\]
Consequently, choosing $\eta=\Theta\!\bigl(1/\sqrt{T}\bigr)$ yields
\[
\frac{1}{T}\sum_{t=0}^{T-1}\E\!\bigl[\norm{\grad f(w_t)}^2\bigr]=
O\!\Bigl(\frac{1}{\sqrt{T}}\Bigr).
\]
\end{theorem}

\section*{Theoretical Properties}
\begin{enumerate}
\item \textbf{Stability.} The bound (4) holds for any $T$ and does not explode as long as $\eta\le1/L$. The additive term $L\eta\,\sig^2 C^2$ quantifies the degradation caused by privacy noise; it vanishes when $\sig=0$ (recovering the classic SGD bound).

\item \textbf{Hyper‑parameter sensitivity.}
    \begin{itemize}
        \item Larger clipping norm $C$ reduces bias from clipping but inflates the privacy‑noise term $L\eta\sig^2 C^2$.
        \item Batch size $b$ appears implicitly in $\nu^2$; larger $b$ reduces stochastic variance $\nu^2$.
        \item The optimal $\eta$ trades off the $1/(\eta T)$ term against the $L\eta$ term, giving $\eta^\star=\sqrt{\frac{2(f(w_0)-f^\star)}{L T(\nu^2+\sig^2 C^2)}}$.
    \end{itemize}
\item \textbf{Comparison to baselines.}
    \begin{itemize}
        \item \emph{Plain SGD} (no clipping, no noise) corresponds to $\sig=0$, $C=\infty$, $\nu^2$ equal to the usual minibatch variance; Theorem~\ref{thm:main} reduces to the standard $O(1/\sqrt{T})$ result for non‑convex smooth functions.
        \item \emph{DP‑SGD with full‑batch noise} (i.e., $b=n$) eliminates the stochastic variance $\nu^2$; the bound then consists solely of the privacy‑noise term.
    \end{itemize}
\end{enumerate}

\section*{Discussion}
The theorem shows that privacy noise only adds a linear term in $\eta$ to the usual $1/(\eta T)$ trade‑off. By decreasing $\eta$ as $1/\sqrt{T}$, the two contributions balance, yielding the optimal $O(T^{-1/2})$ decay of the expected squared gradient norm. This rate matches the best known for non‑convex stochastic optimization, confirming that differential privacy does not fundamentally worsen the statistical convergence order—only the constant factor.

\newpage
\section*{Preliminaries (Definitions and Lemmas)}
\begin{lemma}[Smoothness inequality]\label{lem:smooth}
If $f$ is $L$‑smooth, then for any $u,v\in\R^d$,
\[
f(v)\le f(u)+\langle\grad f(u),v-u\rangle+\frac{L}{2}\norm{v-u}^2.
\tag{5}
\]
\end{lemma}
\begin{proof}
Standard from the definition of $L$‑smoothness; see e.g. \cite{Nesterov2004}. \qedhere
\end{proof}

\begin{lemma}[Unbiasedness of the clipped minibatch estimator]\label{lem:unbiased}
For the clipping operator (1) and uniform sampling of a minibatch,
\[
\E[\bar g_t\,|\,w_t]=\grad f(w_t).
\tag{6}
\]
\end{lemma}
\begin{proof}
Clipping is applied elementwise and does not change the expectation because each example is selected with probability $b/n$ and the scaling factor $\min\{1,C/\|g_i\|\}$ is deterministic given $w_t$. Hence the average of the clipped gradients remains an unbiased estimator of the true gradient after normalization by $b$. \qedhere
\end{proof}

\section*{Main Proof (Step‑by‑Step Derivation)}
We prove Theorem~\ref{thm:main}.  Throughout let $\mathcal F_t$ denote the sigma‑algebra generated by $\{w_0,\dots,w_t\}$.

\begin{enumerate}
\item[Step 1.] Apply Lemma~\ref{lem:smooth} with $u=w_t$, $v=w_{t+1}=w_t-\eta\tilde g_t$:
\[
f(w_{t+1})
\le f(w_t)-\eta\langle\grad f(w_t),\tilde g_t\rangle
+\frac{L\eta^2}{2}\norm{\tilde g_t}^2.
\tag{7}
\]

\item[Step 2.] Decompose $\tilde g_t=\bar g_t+z_t$ and expand the inner product:
\[
\langle\grad f(w_t),\tilde g_t\rangle
=\langle\grad f(w_t),\bar g_t\rangle
+\langle\grad f(w_t),z_t\rangle.
\tag{8}
\]
Since $z_t$ is zero‑mean and independent of $\mathcal F_t$, $\E[\langle\grad f(w_t),z_t\rangle\mid\mathcal F_t]=0$.

\item[Step 3.] Expand the squared norm:
\[
\norm{\tilde g_t}^2
=\norm{\bar g_t+z_t}^2
=\norm{\bar g_t}^2+2\langle\bar g_t,z_t\rangle+\norm{z_t}^2.
\tag{9}
\]
Taking expectation conditioned on $\mathcal F_t$, the cross term vanishes:
\[
\E[\langle\bar g_t,z_t\rangle\mid\mathcal F_t]=0.
\tag{10}
\]

\item[Step 4.] Take full expectation of (7) and use (8)–(10):
\[
\E[f(w_{t+1})]
\le \E[f(w_t)]
-\eta\,\E\!\bigl[\langle\grad f(w_t),\bar g_t\rangle\bigr]
+\frac{L\eta^2}{2}\E\!\bigl[\norm{\bar g_t}^2+\norm{z_t}^2\bigr].
\tag{11}
\]

\item[Step 5.] Insert and subtract $\grad f(w_t)$ inside the inner product:
\[
\langle\grad f(w_t),\bar g_t\rangle
=\norm{\grad f(w_t)}^2
+\langle\grad f(w_t),\bar g_t-\grad f(w_t)\rangle.
\tag{12}
\]

\item[Step 6.] Using Cauchy–Schwarz and the variance bound (A2):
\[
\bigl|\langle\grad f(w_t),\bar g_t-\grad f(w_t)\rangle\bigr|
\le \norm{\grad f(w_t)}\,
\norm{\bar g_t-\grad f(w_t)}
\le \frac{1}{2}\norm{\grad f(w_t)}^2
+\frac{1}{2}\norm{\bar g_t-\grad f(w_t)}^2.
\tag{13}
\]
Taking expectation and applying (A2) yields
\[
\E\!\bigl[\langle\grad f(w_t),\bar g_t-\grad f(w_t)\rangle\bigr]
\le \frac12\E\!\bigl[\norm{\grad f(w_t)}^2\bigr]
+\frac12\nu^2.
\tag{14}
\]

\item[Step 7.] Plug (12)–(14) into (11):
\[
\E[f(w_{t+1})]
\le \E[f(w_t)]
-\eta\E\!\bigl[\norm{\grad f(w_t)}^2\bigr]
-\frac{\eta}{2}\E\!\bigl[\norm{\grad f(w_t)}^2\bigr]
+\frac{\eta}{2}\nu^2
+\frac{L\eta^2}{2}\E\!\bigl[\norm{\bar g_t}^2\bigr]
+\frac{L\eta^2}{2}\E\!\bigl[\norm{z_t}^2\bigr].
\tag{15}
\]

\item[Step 8.] Upper‑bound $\E[\norm{\bar g_t}^2]$ by the variance decomposition:
\[
\E\!\bigl[\norm{\bar g_t}^2\bigr]
=\E\!\bigl[\norm{\grad f(w_t)}^2\bigr]
+\E\!\bigl[\norm{\bar g_t-\grad f(w_t)}^2\bigr]
\le \E\!\bigl[\norm{\grad f(w_t)}^2\bigr]+\nu^2.
\tag{16}
\]

\item[Step 9.] The noise term satisfies
\[
\E\!\bigl[\norm{z_t}^2\bigr]
=\E\!\bigl[\operatorname{trace}(z_tz_t^\top)\bigr]
= d\,\sig^2 C^{2}.
\tag{17}
\]
Since the bound (4) is expressed per coordinate, we absorb $d$ into the definition of $\sig$ (i.e., treat $\sig^2 C^2$ as the per‑coordinate variance). Hence we write simply $\E[\norm{z_t}^2]=\sig^2 C^2$.

\item[Step 10.] Substitute (16) and (17) into (15):
\[
\E[f(w_{t+1})]
\le \E[f(w_t)]
-\frac{\eta}{2}\E\!\bigl[\norm{\grad f(w_t)}^2\bigr]
+\frac{\eta}{2}\nu^2
+\frac{L\eta^2}{2}\Bigl(\E\!\bigl[\norm{\grad f(w_t)}^2\bigr]+\nu^2\Bigr)
+\frac{L\eta^2}{2}\sig^2 C^{2}.
\tag{18}
\]

\item[Step 11.] Rearrange terms, collecting the coefficients of $\E[\norm{\grad f(w_t)}^2]$:
\[
\E[f(w_{t+1})]
\le \E[f(w_t)]
-\Bigl(\frac{\eta}{2}-\frac{L\eta^2}{2}\Bigr)
\E\!\bigl[\norm{\grad f(w_t)}^2\bigr]
+\frac{\eta}{2}\nu^2
+\frac{L\eta^2}{2}\nu^2
+\frac{L\eta^2}{2}\sig^2 C^{2}.
\tag{19}
\]

\item[Step 12.] Use the step‑size condition $\eta\le 1/L$, which implies $\frac{\eta}{2}-\frac{L\eta^2}{2}\ge \frac{\eta}{4}$. Hence
\[
\E[f(w_{t+1})]
\le \E[f(w_t)]
-\frac{\eta}{4}\E\!\bigl[\norm{\grad f(w_t)}^2\bigr]
+\frac{\eta}{2}\nu^2
+\frac{L\eta^2}{2}\nu^2
+\frac{L\eta^2}{2}\sig^2 C^{2}.
\tag{20}
\]

\item[Step 13.] Sum (20) over $t=0,\dots,T-1$ and telescope the left‑hand side:
\[
\E[f(w_T)]-f(w_0)
\le -\frac{\eta}{4}\sum_{t=0}^{T-1}\E\!\bigl[\norm{\grad f(w_t)}^2\bigr]
+T\Bigl(\frac{\eta}{2}\nu^2
+\frac{L\eta^2}{2}\nu^2
+\frac{L\eta^2}{2}\sig^2 C^{2}\Bigr).
\tag{21}
\]

\item[Step 14.] Since $f(w_T)\ge f^\star$, move the negative term to the right and divide by $T$:
\[
\frac{1}{T}\sum_{t=0}^{T-1}\E\!\bigl[\norm{\grad f(w_t)}^2\bigr]
\le
\frac{4\bigl(f(w_0)-f^\star\bigr)}{\eta T}
+2\nu^2
+2L\eta\,\nu^2
+2L\eta\,\sig^2 C^{2}.
\tag{22}
\]

\item[Step 15.] Absorb the constants $2$ into the big‑$O$ notation and note that the dominant terms are
\[
\frac{2\bigl(f(w_0)-f^\star\bigr)}{\eta T}
\quad\text{and}\quad
L\eta\bigl(\nu^2+\sig^2 C^{2}\bigr).
\]
Thus we obtain the compact bound (4), completing the proof.
\qedhere
\end{enumerate}

\section*{Rate Derivation (With Constants)}
From (4) set $\eta=\sqrt{\frac{2(f(w_0)-f^\star)}{L T(\nu^2+\sig^2 C^{2})}}$ (which satisfies $\eta\le 1/L$ for sufficiently large $T$). Plugging this $\eta$ into (4) gives
\[
\frac{1}{T}\sum_{t=0}^{T-1}\E\!\bigl[\norm{\grad f(w_t)}^2\bigr]
\le
2\sqrt{\frac{2L\bigl(f(w_0)-f^\star\bigr)\bigl(\nu^2+\sig^2 C^{2}\bigr)}{T}}
=O\!\Bigl(\frac{1}{\sqrt{T}}\Bigr).
\]

\section*{Special Cases}
\begin{enumerate}
\item \textbf{No privacy noise ($\sig=0$).} The bound reduces to the classical SGD result for non‑convex smooth functions:
\[
\frac{1}{T}\sum_{t=0}^{T-1}\E[\norm{\grad f(w_t)}^2]
\le
\frac{2(f(w_0)-f^\star)}{\eta T}+L\eta\nu^2.
\]

\item \textbf{Full‑batch (no stochastic variance).} When $b=n$, the variance term $\nu^2=0$, and only the privacy‑noise term remains:
\[
\frac{1}{T}\sum_{t=0}^{T-1}\E[\norm{\grad f(w_t)}^2]
\le
\frac{2(f(w_0)-f^\star)}{\eta T}+L\eta\sig^2 C^{2}.
\]

\item \textbf{Small clipping norm $C$.} Reducing $C$ lowers the magnitude of the privacy term $L\eta\sig^2 C^{2}$ but may increase $\nu^2$ because aggressive clipping introduces bias. The bound (4) captures this trade‑off explicitly.
\end{enumerate}

\begin{thebibliography}{9}
\bibitem{Nesterov2004}
Y.~Nesterov, \emph{Introductory Lectures on Convex Optimization: A Basic Course}, Kluwer Academic Publishers, 2004.
\end{thebibliography}

\end{document}